\documentclass[a4paper,11pt]{article}

\usepackage{revy}
\usepackage[utf8]{inputenc}
\usepackage[T1]{fontenc}
\usepackage[danish]{babel}
\usepackage{amsmath,amssymb,amsfonts,amsthm,textcomp}


\revyname{MatematikRevyen}
\revyyear{2026}
\version{0.1}
\eta{$4$ minutter}
\status{Færdig}

\title{PhD Draft}
\author{Louie '24 \& Carl '21}

\begin{document}
\maketitle

\begin{roles}
\role{S1}[Skuespiller] Nyuddannet Kandidat (TalTeori)
\role{S2}[Skuespiller] Nyuddannet Kandidat (OpAlg)
\role{S3}[Skuespiller] Nyuddannet Kandidat (MatHist)
\role{R}[Skuespiller] RUC
\role{K}[Skuespiller] KU
\role{D}[Skuespiller] DTU
\role{S}[Skuespiller] Sportskommentater
\end{roles}


\begin{sketch}

\scene{R, K og D sidder ved et slags dommerbord. S er en hype kommentater, der ikke rigtig ved noget om matematik. }

\says{S} Ve-ve-ve-VELKOMMEN til årets P-H-D. Draaaaft! I kender det! I eeeeelsker det! Det. Her. Det. Er årets store dag, hvor nyudklækkede kandidater finder ud af om de bliver arbejdsløse eller ej! WUuH!

\scene{S finder nogle stikordskort frem og læser op}

\says{S} Anywayz, dagens først kandidat har brugt rigtig rigtig rigtig lang tid på talteo.. Teorien om tal?

\says{S} Jeg troede efterhånden vi havde fået styr på de tal der, men lad os høre hvad de har at sige. Giv den op for S1!

\scene{Publikum klapper og S1 kommer ind på scenen. S1 er lidt en nervøs type}

\says{S1} Hej institutter. Som vi jo selvfølgelig allerede ved, så har Riemann-Zeta funktionen uendelig mange nulpunkter når $\sigma = 0.5$. Men øhh, jeg kunne godt tænke mig.. jeg vil gerne vise, at der er lidt flere 0'er på den linje der. 

\scene{RUC og DTU er begejstrede, men KU er afvisende}

\says{R} Uhhh, fedt!

\says{D} Ja, \textit{det} lyder spændende!

\says{K} Gutter, trust. Jeg ved det lyder flashy med Ri-Ze funktionen, men vi er faldet for den for mange gange før. Et ekstra nulpunkt til eller fra, det gør sku ikke nogen forskel på hverken cardinaliteten eller interessen. 

\scene{S trykker på en 'afvist' buzzer}

\says{S} De-de-denied! Det blev ikke til nogen ph.D denne gang og S1 må desværre vende hjem uden kridt i hånden. Wha-wha

\says{S1}[desperat] Det er altså virkelig spænende! HOV! \act{En ninja hiver S1 af scenen} Nej! Jeg mener vigtigt, vigtiiiig 


\says{S} Tragisk slutning til S1! Lad os håbe den næste har det institutterne søger. 

\scene{S tager stikordkort frem igen og læser højt}

\says{S} Den næste kandidat har specialiseret sig i at.. operere lidt på nogle algebraer? Det lyder anvendt! Giv den op for S2!!

\scene{Publikum klapper og S2 kommer ind på scenen. S2 er en smart type}

\says{S2}[til S] Jeg laver altså Operator Algebra.

\says{S} Yes, Det er fint med dig. Lad os høre hvad institutterne har at sige til din Algebra Opera. RUC?

\says{R} Jeg tror ikke, helt jeg forstår. 

\says{S}[afbryder] Det tænkte jeg sgu nok, der er vi ikke så forskellige os to. Wink. DTU?

\says{D}[tøvende] Ja, altså algebra... Lidt teoretisk, hvad kan dit speciale bruges til? 

\says{K} Ja, det er vi jo normalt ret ligeglade med, meeen man skal jo tænke på fondsansøgninger

\says{S} Det ser desværre ud til....

\scene{S skal til at trykke på 'afvist' buzzer, men S2 afbryder}

\says{S2} Nej nej! vent!! 

\says{S2}[cool] Peep this, hvad nu hvis jeg siger... at det kan bruges til Kante-AI!?! OoOooHhh

\says{S} ÅÅÅHHH, God DamnN! S2 kører den sikkert hjem med et buzzWORD. En fejlfri strategi!

\scene{DTU, KU, RUC er helt begejstrede og kæmper om S2. De rejser sig fra bordet, går op mod S2. Begynder at tilbyde ting}

\says{D} Du SKAL ud til os, vi har...

\says{K} Hjørnekontor

\says{R} Gratis kaffe!

\says{D} Personlig kvantecomputer

\says{R} Parkeringsplads!

\says{K} København!

\scene{S2 river sig løs}

\says{S2} Roooolig babes, der er nok til jer alle sammen. Eller... næsten alle sammen

\scene{S2 tager KU og DTU i hånden og trækker dem frem. RUC er ked af det.}

\scene{S2 lægger armene om KU og DTU}

\says{S2} Hvad med et lille co-lab? Wink

\scene{KU og DTU kysser S2 på kinden og de går alle tre ud sammen}

\scene{S trykker på 'godkendt' buzzer}

\says{S} WOW! Sikke en kamp! Det er sjældent man ser en SÅ attraktiv matematiker. 

\says{S} Det var alt for i år og endnu engang står RUC alene tilbage. \act{RUC er trist} Nååårhgh, måske næste år du

\scene{S får noget at vide i sin øresnegl}

\says{S} Hov vent!! Der er netop i dette øjeblik en studerende, der har afleveret deres speciale. Deadline var jo i sommer, men lad os få dem herind alligevel. Giv den op for S3!!!

\scene{S3 kommer ind}

\says{S} S3 take it away.

\says{S3} Jo, jeg har brugt de sidste 5 år på at kigge historisk på hvorfor folk synes nulpunkter på en linje er spændende. Det er sku lidt et mysterium, så jeg vil gerne have 3 år mere til det.

\says{R}[håbefuld] Vil det sige, at du laver noget \textit{historie}?!? 

\says{S3} Ja, altså. \textit{Matematik} historie

\says{R} Historieeee!! 

\scene{R løber op og krammer S3}

\says{R} Du ville være perfekt til ph.D. i humaniora hos os!

\scene{lys ned.}

\end{sketch}
\end{document}
